\documentclass{situationist-preprint}

\title{The Regress of Why-Questions}
\msid{SCI-03-b}
\msfamily{SCI-03}
\msdiscipline{Philosophy of Science}
\mstype{Consequence}
\mskeywords{explanation; why-questions; regress; principle of sufficient reason; dependence}

\begin{document}
\maketitle

\begin{abstract}
Every explanation cites factors and a generalization, and each of these can be made the topic of a further why-question. It has seemed to follow either that explanation requires an ultimate explanation that ends the regress, or that every explanation rests on something unexplained and so explains nothing. We derive from the account of SCI-03, on which an explanation answers a why-question by locating the explanandum in a pattern of dependence, that neither follows. The regress is real and may be endless, but it is not vicious: an answer is complete relative to its own question. The regress ends locally where a question's presupposition fails or nothing in the relevant range makes a difference, and globally, for science, at the conditions of science itself.
\end{abstract}

\section{The regress}\label{s:regress}

By Principle~P1 of SCI-03-a, an answer to the question why $P$ rather than $P_1, \dots, P_n$ explains if it cites actual values of variables $F_1, \dots, F_k$ and a generalization $G$ that is true of those values, invariant under some range of changes in them, and under some such change assigns a value that makes one of the alternatives true. Each element of the answer is a fact that can be questioned. Why did the $F_j$ take the values they did? Why does $G$ hold, and why does it hold over that range of invariance?

\begin{definition}[Regress]\label{d:regress}
A \emph{regress} of why-questions is a sequence $Q_0, Q_1, \dots$ in which each $Q_{n+1}$ has as its topic a fact cited in an explanation that answers $Q_n$.
\end{definition}

Two inferences have been drawn from the existence of such sequences. The \emph{foundationalist} inference holds that an explanation of $P$ is complete only if each fact it cites is itself explained, and that the regress must therefore terminate in a self-explaining fact, as in the cosmological arguments of the rationalist tradition \citep[§§32--38]{Leibniz1714}. The \emph{skeptical} inference holds that, since the regress cannot terminate in that way, every explanation rests on unexplained facts and so leaves its explanandum as mysterious as it found it.

\section{Consequences of the dependence account}\label{s:consequences}

\begin{proposition}[Locality]\label{p:local}
Whether an answer explains $P$ rather than $P_1, \dots, P_n$ does not depend on whether any fact it cites is explained.
\end{proposition}

\begin{proof}
The three conditions of Principle~P1 of SCI-03-a concern the truth of $G$ of the actual values, its invariance under some range of changes, and what it assigns under such a change. None of them mentions an explanation of the $F_j$ or of $G$. An answer that satisfies them explains, whatever the answers to further questions.
\end{proof}

Proposition~\ref{p:local} defeats both inferences, which share the premise that an explanation inherits its explanatory standing from the explanations of what it cites. On the dependence account, what an explanation conveys is the pattern of dependence of Definition~2.4 of SCI-03-a: the counterfactuals that say how the explanandum would have differed had the factors differed. That information is the same whether or not the factors are themselves explained. The shadow depends on the height of the flagpole in the same way whether or not anyone knows why the pole was built to that height.

\begin{proposition}[Non-termination]\label{p:endless}
Nothing in the account requires a regress to terminate, and nothing requires it to continue.
\end{proposition}

\begin{proof}
By Proposition~\ref{p:local}, an infinite regress does not deprive its members of explanatory standing, so the account does not require termination. Conversely, a regress terminates at $Q_n$ if $Q_n$ has no explanation, and the account allows this in two ways. First, a why-question presupposes that its topic is true and that its alternatives are mutually exclusive \citep{vanFraassen1980}; where the presupposition fails, the question is rejected rather than answered. Second, the question may concern a fact on which no variable within any range of invariance makes a difference, so that no generalization satisfies condition (c) of Principle~P1 of SCI-03-a.
\end{proof}

The second kind of termination is what is usually called a brute fact. On the dependence account a brute fact is not a defect in an explanation that cites it; it is a fact that is not itself the topic of an explanation. The principle of sufficient reason denies that there are such facts \citep{Pruss2006}. The dependence account is neutral on the principle, and it shows that the reason commonly given for it, that otherwise all explanation would be incomplete, rests on the premise Proposition~\ref{p:local} rejects.

\begin{corollary}[The limit of scientific regress]\label{c:science}
Every scientific regress that reaches the existence of the conditions that make science possible terminates as far as science is concerned.
\end{corollary}

\begin{proof}
SCI-18 argues that science cannot explain the existence of the conditions that make it possible, since every scientific explanation cites laws and concrete conditions of the kinds whose existence is in question. A why-question about the existence of all such conditions therefore has no scientific answer. By Proposition~\ref{p:local}, the explanations below it in the regress are not thereby impaired.
\end{proof}

Corollary~\ref{c:science} concerns scientific answers only. SCI-18 leaves open whether such facts have explanations of another kind, and MET-01 offers one for the existence of something rather than nothing.

\section{Objections and replies}\label{s:objections}

\begin{objection}[Understanding]
Even if each answer is complete as an answer, understanding is cumulative: we understand a phenomenon better the further down the regress we can go. So an unexplained explainer leaves understanding incomplete.
\end{objection}

\begin{reply}
We agree that understanding grows as further questions are answered, and nothing in Proposition~\ref{p:local} denies it. What we deny is that understanding of $P$ rather than its alternatives is absent until the regress is completed. Understanding comes in degrees, and each answered question in a regress adds to it \citep{Garfinkel1981}.
\end{reply}

\begin{objection}[Generalizations]
The facts cited are of two kinds. Particular values may be explained by earlier values, but a generalization that is explained by nothing is a mere regularity, and a mere regularity cannot explain.
\end{objection}

\begin{reply}
The objection confuses being unexplained with being accidental. What gives a generalization its explanatory standing on the account of SCI-03-a is invariance under changes in its variables \citep{Woodward2003}, and a generalization can be invariant whether or not a deeper theory explains why. The fundamental laws, if there are any, are invariant and unexplained.
\end{reply}

\section{Conclusion}\label{s:conclusion}

On the account of SCI-03, each explanation is complete relative to its question, and its standing is not borrowed from the explanation of what it cites. The regress of why-questions is therefore harmless whether it is infinite or terminates. It terminates locally at rejected questions and brute facts, and globally, for science, at the conditions that make science possible.

\begin{thebibliography}{9}
\bibitem[Garfinkel(1981)]{Garfinkel1981} Garfinkel, A. (1981). \emph{Forms of Explanation: Rethinking the Questions in Social Theory}. Yale University Press.
\bibitem[Leibniz(1714)]{Leibniz1714} Leibniz, G.~W. (1714). \emph{Monadology}. In R.~Ariew and D.~Garber (eds. and trans.), \emph{Philosophical Essays}, Hackett, 1989.
\bibitem[Pruss(2006)]{Pruss2006} Pruss, A.~R. (2006). \emph{The Principle of Sufficient Reason: A Reassessment}. Cambridge University Press.
\bibitem[van Fraassen(1980)]{vanFraassen1980} van Fraassen, B.~C. (1980). \emph{The Scientific Image}. Clarendon Press.
\bibitem[Woodward(2003)]{Woodward2003} Woodward, J. (2003). \emph{Making Things Happen: A Theory of Causal Explanation}. Oxford University Press.
\end{thebibliography}

\end{document}
